\documentclass[10pt,twocolumn]{article}

\usepackage[margin=1in]{geometry}
\usepackage{amsmath,amssymb}
\usepackage{graphicx}
\usepackage{booktabs}
\usepackage{hyperref}
\usepackage{xcolor}
\usepackage{microtype}

\hypersetup{
  colorlinks=true,
  linkcolor=blue!70!black,
  citecolor=blue!70!black,
  urlcolor=blue!70!black,
}

% -----------------------------------------------------------------------
\title{\textbf{CTDenoiser}: A Benchmark for Low-Dose CT Denoising
       Across Supervised and Self-Supervised Regimes}

% -----------------------------------------------------------------------

\begin{document}
\maketitle

% -----------------------------------------------------------------------
\begin{abstract}
Low-dose computed tomography (LDCT) denoising must contend with
structured, spatially correlated noise that ordinary blind-spot
self-supervision is not designed to handle.  We present \textbf{CTDenoiser},
a benchmark and single data/training/inference harness that compares
denoisers across five architectures (RED-CNN, DnCNN, U-Net, CTformer, and a
conditional flow-matching model, CFM) and five supervision regimes:
fully supervised training with paired full-dose targets, Noise2Sim
similarity-based self-supervision, Noise2Void blind-spot self-supervision,
Zero-Shot Noise2Noise (ZS-N2N), and Filter2Noise.  Models are scored with
PSNR, SSIM, and RMSE alongside Gradient Magnitude Similarity Deviation
(GMSD) and a Noise Power Spectrum ratio (NPS-ratio) that targets spurious
high-frequency texture.  On the Mayo TCIA abdomen data, supervised models
recover $\sim$4\,dB PSNR and $+0.046$ SSIM over the noisy input and finish
within $0.1$\,dB of one another across CNN backbones, indicating that the
supervision signal matters far more than the architecture.  Among methods
that use no clean targets, only the correlated-noise-aware Noise2Sim
consistently beats the noisy baseline ($+0.5$ to $+1.2$\,dB), while
blind-spot Noise2Void is flat-to-negative (down to $-0.46$\,dB), direct
evidence that CT's spatially correlated noise violates the
pixel-independence assumption behind blind-spot and downsampling-based
methods.  The resulting $3$--$4$\,dB gap between self-supervised and
supervised denoising motivates correlated-noise-aware self-supervision as
the key open problem.  Results are single-seed on abdomen at 50 epochs;
multi-anatomy and multi-seed runs are in progress.
\end{abstract}

% -----------------------------------------------------------------------
\section{Related Work}

\paragraph{CNN-based denoisers.}
RED-CNN~\cite{chen2017lowdose} introduced a residual encoder-decoder with
multi-level skip connections and remains a standard baseline.
DnCNN~\cite{zhang2017beyond} showed that a flat stack of Conv-BN-ReLU
layers predicting the residual (noise) transfers well across denoising
tasks.

\paragraph{Transformer and generative denoisers.}
CTformer~\cite{wang2022ctformer} applies a U-shaped transformer with
Token2Token dilation to capture long-range structure at reduced memory
cost.  Score-based diffusion~\cite{song2021scorebased,gao2023corediff}
can yield perceptually strong results but needs many neural function
evaluations; conditional flow matching~\cite{lipman2022flow,albergo2022building}
learns near-straight transport paths and cuts this to tens of steps.

\paragraph{Self-supervised and zero-shot denoisers.}
When paired full-dose targets are unavailable, denoisers can train from
noisy data alone.  Noise2Noise~\cite{lehtinen2018noise2noise} maps between
two noisy observations of a scene; Noise2Void~\cite{krull2019noise2void}
removes the second observation via a blind-spot scheme but assumes
pixel-wise independent noise, an assumption CT violates.  Two lines
of work target correlated noise directly:
Noise2Sim~\cite{niu2020noise2sim} maps between centre pixels of
self-similar non-local patches, and Zero-Shot Noise2Noise
(ZS-N2N)~\cite{mansour2023zsn2n} trains a tiny per-image network from a
downsampled Noise2Noise pair with no external data.  More recently,
Filter2Noise~\cite{sun2025filter2noise} couples single-image
self-supervision with an interpretable attention-guided bilateral filter
for LDCT.  Our harness includes Noise2Void and ZS-N2N as blind-spot and
zero-shot baselines and adds Noise2Sim and Filter2Noise as the
correlated-noise-aware extensions.

\paragraph{Evaluation metrics.}
PSNR and SSIM~\cite{wang2004image} favour pixel-level fidelity and can
penalise plausible-but-incorrect high-frequency detail.  We complement
them with GMSD~\cite{xue2014gradient} as an efficient perceptual proxy and
an NPS-ratio inspired by noise-texture
characterisation~\cite{solomon2012characterization} to flag spurious
residual texture; LPIPS~\cite{zhang2018unreasonable} is a planned addition.

% -----------------------------------------------------------------------
\section{Experimental Setup}
We train on the Mayo TCIA LDCT-and-Projection-Data~\cite{moen2021low}
(abdomen window), with a patient-level split (seed 0): $6{,}683$ training
and $1{,}476$ validation slices, HU windowed to $[0,1]$.  Trainable models
run for 50 epochs; the per-image methods (ZS-N2N, Filter2Noise) optimise at
test time and are model-agnostic.  Validation metrics are computed on full
slices via overlapped tiling; we corrected a border-handling bug in that
tiling that previously zeroed each slice's outer margin and biased
reconstruction metrics against the trained models.  We report the gain
($\Delta$) of each metric relative to the \emph{LDCT input} (the noisy scan
with no denoising), which is the floor every method must beat.

% -----------------------------------------------------------------------
\section{Results}

\begin{table*}[t]
  \centering
  \small
  \caption{Validation results on TCIA LDCT (abdomen; patient-level split,
           seed 0; 50 epochs), grouped by supervision regime.  Means over
           the validation set on full slices via overlapped tiling
           (corrected border handling).  PSNR shows the gain over the
           \emph{LDCT input} floor in parentheses.  $\uparrow$/$\downarrow$:
           higher/lower is better.  \textbf{Bold}: best denoiser per metric.
           ZS-N2N and Filter2Noise are per-image and model-agnostic (single
           row each).  $\dagger$: run truncated by a crash, reported at its
           last logged epoch (CTformer/Supervised ep.\,37, CTformer/Noise2Sim
           ep.\,26).  $\ddagger$: CFM crashed at ep.\,15 and had not
           converged; shown for reference.  CFM requires paired targets, so
           it has no blind-spot or similarity variant.}
  \label{tab:results}
  \begin{tabular}{llrlcccc}
    \toprule
    \textbf{Model} & \textbf{Supervision} & \textbf{Params} &
    PSNR\,$\uparrow$ (\,$\Delta$) & SSIM\,$\uparrow$ & RMSE\,$\downarrow$ &
    GMSD\,$\downarrow$ & NPS-r.\,$\downarrow$ \\
    \midrule
    LDCT input & --- (no denoising) & --- & 31.44 & 0.893 & 0.0274 & 0.166 & 0.0079 \\
    \midrule
    RED-CNN~\cite{chen2017lowdose}   & Supervised & $1.85\mathrm{M}$ & \textbf{35.48} \,(+4.04) & \textbf{0.940} & 0.0171 & 0.166 & 0.0031 \\
    U-Net                            & Supervised & $1.95\mathrm{M}$ & 35.47 \,(+4.04) & 0.940 & \textbf{0.0171} & 0.178 & \textbf{0.0030} \\
    DnCNN~\cite{zhang2017beyond}     & Supervised & $0.56\mathrm{M}$ & 35.40 \,(+3.97) & 0.938 & 0.0172 & 0.184 & 0.0031 \\
    CTformer~\cite{wang2022ctformer}$^{\dagger}$ & Supervised & $0.40\mathrm{M}$ & 34.78 \,(+3.34) & 0.917 & 0.0184 & 0.263 & 0.0036 \\
    CFM~\cite{lipman2022flow}$^{\ddagger}$       & Supervised & $2.11\mathrm{M}$ & 33.87 \,(+2.43) & 0.923 & 0.0206 & 0.248 & 0.0044 \\
    \midrule
    U-Net                            & Noise2Sim~\cite{niu2020noise2sim} & $1.95\mathrm{M}$ & 32.61 \,(+1.18) & 0.910 & 0.0238 & 0.169 & 0.0060 \\
    DnCNN~\cite{zhang2017beyond}     & Noise2Sim~\cite{niu2020noise2sim} & $0.56\mathrm{M}$ & 32.14 \,(+0.70) & 0.891 & 0.0251 & 0.178 & 0.0067 \\
    RED-CNN~\cite{chen2017lowdose}   & Noise2Sim~\cite{niu2020noise2sim} & $1.85\mathrm{M}$ & 31.94 \,(+0.51) & 0.904 & 0.0257 & \textbf{0.162} & 0.0070 \\
    CTformer~\cite{wang2022ctformer}$^{\dagger}$ & Noise2Sim~\cite{niu2020noise2sim} & $0.40\mathrm{M}$ & 31.09 \,($-0.35$) & 0.885 & 0.0282 & 0.247 & 0.0083 \\
    \midrule
    RED-CNN~\cite{chen2017lowdose}   & Noise2Void~\cite{krull2019noise2void} & $1.85\mathrm{M}$ & 31.52 \,(+0.08) & 0.895 & 0.0271 & 0.175 & 0.0077 \\
    U-Net                            & Noise2Void~\cite{krull2019noise2void} & $1.95\mathrm{M}$ & 31.49 \,(+0.05) & 0.893 & 0.0271 & 0.188 & 0.0078 \\
    CTformer~\cite{wang2022ctformer} & Noise2Void~\cite{krull2019noise2void} & $0.40\mathrm{M}$ & 31.28 \,($-0.15$) & 0.894 & 0.0277 & 0.223 & 0.0081 \\
    DnCNN~\cite{zhang2017beyond}     & Noise2Void~\cite{krull2019noise2void} & $0.56\mathrm{M}$ & 30.98 \,($-0.46$) & 0.885 & 0.0288 & 0.198 & 0.0087 \\
    \midrule
    Filter2Noise~\cite{sun2025filter2noise} & Zero-shot & --- & 31.58 \,(+0.14) & 0.896 & 0.0269 & 0.176 & 0.0076 \\
    ZS-N2N~\cite{mansour2023zsn2n}          & Zero-shot & $21.3\mathrm{K}$ & 31.17 \,($-0.27$) & 0.892 & 0.0281 & 0.182 & 0.0083 \\
    \bottomrule
  \end{tabular}
\end{table*}

Table~\ref{tab:results} yields three findings.  \textbf{First}, supervised
denoising gains $\sim$4\,dB PSNR and $+0.046$ SSIM over the noisy input, and
RED-CNN, U-Net, and DnCNN finish within $0.1$\,dB of one another despite a
$3\times$ range in parameter count, so the objective dominates the backbone.
\textbf{Second}, among label-free methods only the correlated-noise-aware
Noise2Sim consistently exceeds the noisy floor ($+0.5$ to $+1.2$\,dB),
whereas blind-spot Noise2Void is flat-to-negative ($-0.46$\,dB and $-0.008$
SSIM on DnCNN) and the zero-shot ZS-N2N sits essentially at the floor; this
is direct evidence that FBP-correlated CT noise violates the
pixel-independence assumption these methods rely on.  \textbf{Third}, the
$3$--$4$\,dB gap between the best self-supervised result and the supervised
upper bound quantifies the headroom for correlated-noise-aware
self-supervision.  CTformer and CFM runs were truncated by crashes and are
reported at their last logged epoch; completing them, together with
multi-anatomy and multi-seed evaluation, is in progress.

% -----------------------------------------------------------------------
\bibliographystyle{plain}
\begin{thebibliography}{99}

\bibitem{brenner2007computed}
D.~J. Brenner and E.~J. Hall.
\newblock Computed tomography, an increasing source of radiation exposure.
\newblock \emph{New England Journal of Medicine}, 357(22):2277--2284, 2007.

\bibitem{chen2017lowdose}
H.~Chen, Y.~Zhang, M.~K. Kalra, F.~Lin, Y.~Chen, P.~Liao, J.~Zhou, and
  G.~Wang.
\newblock Low-dose {CT} with a residual encoder-decoder convolutional neural
  network.
\newblock \emph{IEEE Transactions on Medical Imaging}, 36(12):2524--2535, 2017.

\bibitem{zhang2017beyond}
K.~Zhang, W.~Zuo, Y.~Chen, D.~Meng, and L.~Zhang.
\newblock Beyond a {G}aussian denoiser: Residual learning of deep {CNN} for
  image denoising.
\newblock \emph{IEEE Transactions on Image Processing}, 26(7):3142--3155, 2017.

\bibitem{wang2022ctformer}
C.~Wang, S.~Shang, H.~Zhang, Q.~Li, and S.~K. Zhou.
\newblock {CTformer}: Convolution-free token2token dilated faster self-attention
  for {CT} image denoising.
\newblock \emph{Physics in Medicine \& Biology}, 68(6):065012, 2023.

\bibitem{moen2021low}
T.~R. Moen, B.~Chen, D.~R. Holmes, X.~Duan, Z.~Yu, L.~Yu, S.~Leng,
  J.~G. Fletcher, and C.~H. McCollough.
\newblock Low-dose {CT} image and projection data.
\newblock \emph{Medical Physics}, 48(2):902--911, 2021.

\bibitem{song2021scorebased}
Y.~Song, J.~Sohl-Dickstein, D.~P. Kingma, A.~Kumar, S.~Ermon, and B.~Poole.
\newblock Score-based generative modeling through stochastic differential
  equations.
\newblock In \emph{International Conference on Learning Representations}, 2021.

\bibitem{gao2023corediff}
Q.~Gao, F.~Li, Y.~Zhang, Y.~Wang, X.~Hu, P.~Liu, W.~Zheng, H.~Chen, and
  S.~Li.
\newblock {CoreDiff}: Contextual error-modulated generalized diffusion model for
  low-dose {CT} denoising and generalization.
\newblock \emph{IEEE Transactions on Medical Imaging}, 43(2):745--759, 2024.

\bibitem{lipman2022flow}
Y.~Lipman, R.~T.~Q. Chen, H.~Ben-Hamu, M.~Nickel, and M.~Le.
\newblock Flow matching for generative modeling.
\newblock In \emph{International Conference on Learning Representations}, 2023.

\bibitem{albergo2022building}
M.~S. Albergo and E.~Vanden-Eijnden.
\newblock Building normalizing flows with stochastic interpolants.
\newblock In \emph{International Conference on Learning Representations}, 2023.

\bibitem{lehtinen2018noise2noise}
J.~Lehtinen, J.~Munkberg, J.~Hasselgren, S.~Laine, T.~Karras, M.~Aittala,
  and T.~Aila.
\newblock {Noise2Noise}: Learning image restoration without clean data.
\newblock In \emph{International Conference on Machine Learning}, 2018.

\bibitem{krull2019noise2void}
A.~Krull, T.-O. Buchholz, and F.~Jug.
\newblock {Noise2Void}, learning denoising from single noisy images.
\newblock In \emph{IEEE Conference on Computer Vision and Pattern Recognition},
  2019.

\bibitem{niu2020noise2sim}
C.~Niu and G.~Wang.
\newblock {Noise2Sim}, similarity-based self-learning for image denoising.
\newblock \emph{arXiv preprint arXiv:2011.03384}, 2020.

\bibitem{mansour2023zsn2n}
Y.~Mansour and R.~Heckel.
\newblock Zero-shot {Noise2Noise}: Efficient image denoising without any data.
\newblock In \emph{IEEE Conference on Computer Vision and Pattern Recognition},
  2023.

\bibitem{sun2025filter2noise}
Y.~Sun, L.-S. Schneider, M.~Gu, S.~Mei, C.~Ye, F.~Wagner, S.~Bayer, and
  A.~Maier.
\newblock {Filter2Noise}: Interpretable self-supervised single-image denoising
  for low-dose {CT} with attention-guided bilateral filtering.
\newblock \emph{arXiv preprint arXiv:2504.13519}, 2025.

\bibitem{wang2004image}
Z.~Wang, A.~C. Bovik, H.~R. Sheikh, and E.~P. Simoncelli.
\newblock Image quality assessment: From error visibility to structural
  similarity.
\newblock \emph{IEEE Transactions on Image Processing}, 13(4):600--612, 2004.

\bibitem{xue2014gradient}
W.~Xue, L.~Zhang, X.~Mou, and A.~C. Bovik.
\newblock Gradient magnitude similarity deviation: A highly efficient perceptual
  image quality index.
\newblock \emph{IEEE Transactions on Image Processing}, 23(2):684--695, 2014.

\bibitem{zhang2018unreasonable}
R.~Zhang, P.~Isola, A.~A. Efros, E.~Shechtman, and O.~Wang.
\newblock The unreasonable effectiveness of deep features as a perceptual
  metric.
\newblock In \emph{IEEE Conference on Computer Vision and Pattern Recognition},
  2018.

\bibitem{solomon2012characterization}
J.~B. Solomon, O.~Christianson, and E.~Samei.
\newblock Quantitative comparison of noise texture across {CT} scanners from
  different manufacturers.
\newblock \emph{Medical Physics}, 39(10):6048--6055, 2012.

\end{thebibliography}

\end{document}
